\documentclass[12pt,a4paper]{article}

% ---------------------------------------------------------------------------
%  Encodage, langue, polices
% ---------------------------------------------------------------------------
\usepackage[utf8]{inputenc}
\usepackage[T1]{fontenc}
\usepackage[french]{babel}
\usepackage{lmodern}

% ---------------------------------------------------------------------------
%  Mise en page
% ---------------------------------------------------------------------------
\usepackage[top=2.5cm, bottom=2.5cm, left=2.8cm, right=2.8cm]{geometry}
\usepackage{parskip}          % espace entre paragraphes, pas d'indentation
\setlength{\parskip}{6pt}

% ---------------------------------------------------------------------------
%  Couleurs et mise en forme
% ---------------------------------------------------------------------------
\usepackage{xcolor}
\definecolor{kafkared}{RGB}{211,47,47}
\definecolor{flinkblue}{RGB}{21,101,192}
\definecolor{airflowgreen}{RGB}{30,130,76}
\definecolor{codebackground}{RGB}{245,245,245}
\definecolor{bordergray}{RGB}{200,200,200}
\definecolor{warnorange}{RGB}{230,81,0}
\definecolor{titlegray}{RGB}{55,55,55}
\definecolor{pgblue}{RGB}{0,60,120}

% ---------------------------------------------------------------------------
%  Listings (code source)
% ---------------------------------------------------------------------------
\usepackage{listings}
\usepackage{lstautogobble}

\lstdefinestyle{bash}{
  language=bash,
  basicstyle=\ttfamily\small,
  backgroundcolor=\color{codebackground},
  frame=single,
  rulecolor=\color{bordergray},
  breaklines=true,
  autogobble=true,
  showstringspaces=false,
  keywordstyle=\color{flinkblue}\bfseries,
  commentstyle=\color{gray}\itshape,
}

\lstdefinestyle{python}{
  language=Python,
  basicstyle=\ttfamily\small,
  backgroundcolor=\color{codebackground},
  frame=single,
  rulecolor=\color{bordergray},
  breaklines=true,
  autogobble=true,
  showstringspaces=false,
  keywordstyle=\color{flinkblue}\bfseries,
  stringstyle=\color{kafkared},
  commentstyle=\color{gray}\itshape,
}

\lstdefinestyle{sql}{
  language=SQL,
  basicstyle=\ttfamily\small,
  backgroundcolor=\color{codebackground},
  frame=single,
  rulecolor=\color{bordergray},
  breaklines=true,
  autogobble=true,
  showstringspaces=false,
  extendedchars=true,
  inputencoding=utf8,
  keywordstyle=\color{flinkblue}\bfseries,
  stringstyle=\color{kafkared},
  commentstyle=\color{gray}\itshape,
  morekeywords={PROCTIME,CURRENT_TIMESTAMP,IF,NOT,EXISTS,STRING,DOUBLE,BIGINT},
}

\lstdefinestyle{json}{
  basicstyle=\ttfamily\scriptsize,
  backgroundcolor=\color{codebackground},
  frame=single,
  rulecolor=\color{bordergray},
  breaklines=true,
  showstringspaces=false,
}

% ---------------------------------------------------------------------------
%  Encadrés (boîtes d'info / attention / concept)
% ---------------------------------------------------------------------------
\usepackage{tcolorbox}
\tcbuselibrary{skins, breakable}

\newtcolorbox{conceptbox}[1]{
  colback=flinkblue!5,
  colframe=flinkblue!70,
  fonttitle=\bfseries\small,
  title={Concept clé --- #1},
  breakable,
}

\newtcolorbox{warningbox}{
  colback=warnorange!8,
  colframe=warnorange,
  fonttitle=\bfseries\small,
  title={Attention},
  breakable,
}

\newtcolorbox{infobox}[1]{
  colback=airflowgreen!5,
  colframe=airflowgreen!60,
  fonttitle=\bfseries\small,
  title={#1},
  breakable,
}

\newtcolorbox{fournibox}{
  colback=pgblue!5,
  colframe=pgblue!60,
  fonttitle=\bfseries\small,
  title={Fichier fourni --- à lire, pas à modifier},
  breakable,
}

% ---------------------------------------------------------------------------
%  Tableaux
% ---------------------------------------------------------------------------
\usepackage{booktabs}
\usepackage{array}
\usepackage{tabularx}

% ---------------------------------------------------------------------------
%  Liens, méta-données PDF
% ---------------------------------------------------------------------------
\usepackage{hyperref}
\hypersetup{
  colorlinks=true,
  linkcolor=flinkblue,
  urlcolor=kafkared,
  citecolor=gray,
  pdftitle={TP Data Engineering Temps Réel --- Pipeline Séismes},
  pdfauthor={Golden Collar Academy},
}

% ---------------------------------------------------------------------------
%  En-têtes / pieds de page
% ---------------------------------------------------------------------------
\usepackage{fancyhdr}
\pagestyle{fancy}
\fancyhf{}
\fancyhead[L]{\small\textcolor{titlegray}{Data Engineering Temps Réel --- Séance 7}}
\fancyhead[R]{\small\textcolor{titlegray}{TP Séismes}}
\fancyfoot[C]{\small\textcolor{gray}{\thepage}}
\renewcommand{\headrulewidth}{0.4pt}

% ---------------------------------------------------------------------------
%  Numérotation des sections
% ---------------------------------------------------------------------------
\usepackage{titlesec}
\titleformat{\section}{\large\bfseries\color{titlegray}}{Partie \thesection.}{0.5em}{}
\titleformat{\subsection}{\normalsize\bfseries\color{titlegray}}{\thesubsection.}{0.5em}{}

% ---------------------------------------------------------------------------
%  Divers
% ---------------------------------------------------------------------------
\usepackage{enumitem}
\usepackage{graphicx}
\usepackage{amsmath}

% ===========================================================================
%  DÉBUT DU DOCUMENT
% ===========================================================================
\begin{document}

% ---------------------------------------------------------------------------
%  PAGE DE TITRE
% ---------------------------------------------------------------------------
\begin{titlepage}
  \centering
  \vspace*{2cm}

  {\LARGE\bfseries\color{titlegray} Data Engineering Temps Réel}\\[0.4cm]
  {\large\color{gray} Séance 7 --- Travaux Pratiques}\\[2cm]

  {\Huge\bfseries\color{kafkared} Pipeline Séismes}\\[0.3cm]
  {\Large\color{flinkblue} USGS API $\to$ Kafka $\to$ Flink SQL $\to$ PostgreSQL $\to$ Tableau}\\[2.5cm]

  \begin{tabular}{ll}
    \textbf{Durée estimée} & 3 heures \\[0.3cm]
    \textbf{Niveau}        & BAC+5 --- Data Engineering \\[0.3cm]
    \textbf{Prérequis}     & Séances 1--6 (Kafka, Flink SQL, Airflow) \\[0.3cm]
    \textbf{Technologies}  & Python, Kafka, Flink SQL, Airflow, PostgreSQL, Tableau \\
  \end{tabular}

  \vfill
  {\small\color{gray} Golden Collar Academy --- \today}
\end{titlepage}

% ---------------------------------------------------------------------------
%  TABLE DES MATIÈRES
% ---------------------------------------------------------------------------

% ===========================================================================
%  INTRODUCTION
% ===========================================================================
\section*{Introduction et objectifs}
\addcontentsline{toc}{section}{Introduction et objectifs}

\subsection*{Contexte}

Les séismes sont l'un des phénomènes naturels les plus dangereux et les plus imprévisibles.
Chaque jour, des centaines de tremblements de terre sont enregistrés dans le monde entier
par un réseau mondial de sismographes.

Le \textbf{U.S. Geological Survey (USGS)} centralise ces données et les expose via une
API publique en temps réel. Cette API est idéale pour un TP de Data Engineering :
elle est gratuite, sans clé d'API, et produit des données réelles en continu.

\subsection*{Ce que vous recevez}

Pour ce TP, vous disposez de deux dossiers déjà implémentés :

\begin{itemize}
  \item \texttt{infra/} : le fichier \texttt{docker-compose.yml} qui démarre l'ensemble
        des services (Kafka, Flink, Airflow, PostgreSQL, Kafdrop).
  \item \texttt{producer/} : le script Python \texttt{earthquake\_producer.py} qui
        interroge l'API USGS toutes les 30 secondes et publie les séismes dans Kafka
        \textbf{et} PostgreSQL en temps réel.
\end{itemize}

\subsection*{Ce que vous devez construire}

\begin{enumerate}[label=\textbf{\arabic*.}]
  \item \textbf{Le DAG Airflow} (\texttt{dags/earthquake\_pipeline.py}) :
        orchestre le démarrage du pipeline --- crée les topics Kafka,
        soumet les jobs Flink, vérifie que tout fonctionne.
  \item \textbf{Le fichier Flink SQL} (\texttt{flink-jobs/earthquake\_processing.sql}) :
        enrichit chaque séisme en temps réel (catégorie de profondeur, sévérité)
        et génère des alertes pour les séismes dangereux.
\end{enumerate}

\subsection*{Objectifs du TP}

À la fin de ce TP, vous serez capables de :

\begin{enumerate}[label=\textbf{\arabic*.}]
  \item Comprendre l'architecture d'un producteur temps réel (polling + Kafka + PostgreSQL)
  \item Créer un DAG Airflow d'initialisation de pipeline
  \item Écrire des jobs Flink SQL d'enrichissement \textbf{sans fenêtre temporelle}
  \item Vérifier le flux de bout en bout via Kafdrop et l'interface Flink
  \item Visualiser les données sismiques dans Tableau Desktop via PostgreSQL
\end{enumerate}

\subsection*{Architecture du pipeline}

\begin{center}
\begin{tabular}{ccccccccc}
  \textbf{\color{pgblue}Producer} &
  $\to$ &
  \textbf{\color{kafkared}Kafka} &
  $\to$ &
  \textbf{\color{flinkblue}Flink SQL} &
  $\to$ &
  \textbf{\color{kafkared}Kafka} &
  $\to$ &
  \textbf{Tableau} \\[0.2cm]
  \small(fourni) & &
  \small\texttt{earthquake-raw} & &
  \small(à construire) & &
  \small\texttt{earthquake-enriched} & & \\[0.3cm]
  & & & & & & \textbf{\color{pgblue}PostgreSQL} & $\to$ & \textbf{Tableau} \\
  & & & & & & \small\texttt{earthquake\_events} & & \\
\end{tabular}
\end{center}

\vspace{0.5cm}

\begin{infobox}{Rôles des composants}
  \begin{itemize}
    \item \textbf{Producer (Docker, permanent)} : interroge USGS toutes les 30s,
          écrit dans \texttt{earthquake-raw} (Kafka) et \texttt{earthquake\_events} (PostgreSQL).
    \item \textbf{kafka-to-postgres (Docker, permanent)} : consomme les topics Flink enrichis
          (\texttt{earthquake-alerts}, etc.) et les matérialise dans la base \texttt{earthquake} (PostgreSQL).
    \item \textbf{Airflow DAG} : déclenché \emph{toutes les 30 minutes},
          vérifie l'infrastructure, crée les topics, soumet les jobs Flink, vérifie que le sink tourne.
    \item \textbf{Flink SQL} : lit \texttt{earthquake-raw} en continu,
          enrichit chaque événement et publie dans \texttt{earthquake-enriched} ;
          filtre les alertes dans \texttt{earthquake-alerts}.
    \item \textbf{Tableau Desktop} : se connecte à PostgreSQL
          (base \texttt{earthquake}) pour afficher la carte mondiale et les alertes.
  \end{itemize}
\end{infobox}

Les topics Kafka du pipeline :

\begin{center}
\begin{tabularx}{\textwidth}{lX}
  \toprule
  \textbf{Topic Kafka} & \textbf{Contenu} \\
  \midrule
  \texttt{earthquake-raw}      & Données brutes USGS (publiées par le producer) \\
  \texttt{earthquake-enriched} & Séismes enrichis par Flink (profondeur, sévérité, significativité) \\
  \texttt{earthquake-alerts}   & Alertes pour séismes $\geq$ 5.0 ou alerte tsunami \\
  \bottomrule
\end{tabularx}
\end{center}

\newpage

% ===========================================================================
%  PARTIE 1 — EXPLORATION DE L'API USGS
% ===========================================================================
\section{Exploration de l'API USGS}

\subsection{L'API USGS Earthquake}

L'API USGS FDSN (Federated Data System Network) expose les données sismiques mondiales.
Elle retourne du \textbf{GeoJSON} (RFC 7946), un format standard pour les données géospatiales.

\begin{infobox}{Documentation de référence}
  \url{https://earthquake.usgs.gov/fdsnws/event/1/}
\end{infobox}

\subsection{Structure de la réponse GeoJSON}

Un appel à l'API retourne une \textit{FeatureCollection} GeoJSON.
Chaque \texttt{feature} représente un séisme :

\begin{lstlisting}[style=json, caption={Structure d'un séisme GeoJSON}]
{
  "type": "Feature",
  "properties": {
    "mag":      3.0,
    "place":    "22 km WNW of Mentone, Texas",
    "time":     1774913909371,   // millisecondes epoch UTC
    "status":   "reviewed",
    "tsunami":  0,
    "sig":      138,             // indice de significativite (0-1000+)
    "magType":  "ml",
    "type":     "earthquake",
    "title":    "M 3.0 - 22 km WNW of Mentone, Texas"
  },
  "geometry": {
    "type": "Point",
    "coordinates": [-103.832, 31.757, 7.6]  // [longitude, latitude, profondeur_km]
  },
  "id": "tx2026ghixus"
}
\end{lstlisting}

\begin{warningbox}
  \textbf{Attention à l'ordre des coordonnées GeoJSON :}
  \texttt{coordinates} = \texttt{[longitude, latitude, profondeur]}.
  C'est l'inverse de la convention habituelle (lat, lon) !
\end{warningbox}

\subsection{Test de l'API en ligne de commande}

\begin{lstlisting}[style=bash, caption={Appel curl vers l'API USGS}]
# Recuperer les seismes des dernières 24h de magnitude >= 2.5
$start = (Get-Date).AddDays(-1).ToUniversalTime().ToString("yyyy-MM-ddTHH:mm:ss")
$end = (Get-Date).ToUniversalTime().ToString("yyyy-MM-ddTHH:mm:ss")
$url = "https://earthquake.usgs.gov/fdsnws/event/1/query?format=geojson&starttime=$start&endtime=$end&minmagnitude=2.5&limit=5"

curl.exe -s $url | python -m json.tool | Select-Object -First 80
\end{lstlisting}

\subsection{Analyse du producer fourni}

Ouvrez le fichier \texttt{producer/earthquake\_producer.py}.

\begin{fournibox}
  Ce fichier vous est fourni. Il tourne comme un service Docker permanent.
  Vous n'avez pas à le modifier, mais comprenez-le : il alimente tout le pipeline.
\end{fournibox}

Repérez les constantes de configuration en haut du fichier :

\begin{lstlisting}[style=python, caption={Constantes du producer}]
POLL_INTERVAL         = 30   # secondes entre chaque appel USGS
POLL_OVERLAP_SEC      = 90   # chevauchement entre polls (evite les trous)
INITIAL_LOOKBACK_DAYS = 7    # premier poll : remonte N jours pour amorcer
\end{lstlisting}

\begin{conceptbox}{Seeding initial vs polls suivants}
  Le premier poll remonte \textbf{7 jours} d'historique, ce qui garantit d'avoir
  immédiatement des données dans Kafka et PostgreSQL même si peu de séismes
  se sont produits ces dernières minutes.

  Les polls suivants utilisent un chevauchement de 90 secondes (\texttt{POLL\_OVERLAP\_SEC})
  pour ne pas rater d'événements publiés tardivement par l'API USGS.
\end{conceptbox}

\begin{warningbox}
  \textbf{Pourquoi le topic earthquake-raw semble "figé" après le seeding initial ?}

  Le producer tourne correctement et interroge l'API toutes les 30s ---
  vous pouvez le vérifier avec \texttt{docker logs earthquake-producer}.
  Mais chaque poll ne regarde que les 90 dernières secondes (\texttt{POLL\_OVERLAP\_SEC}).

  \medskip
  Globalement, un séisme de magnitude $\geq$ 2.5 se produit environ
  \textbf{toutes les 15 à 30 minutes}. La probabilité d'en avoir un
  dans une fenêtre de 90 secondes est d'environ 10\,\%.
  Il est donc normal d'observer 20 à 30 minutes sans aucun nouvel événement.

  \medskip
  \textbf{Si vous souhaitez un flux plus dense pour les tests}, vous pouvez
  temporairement baisser \texttt{MIN\_MAGNITUDE} à \texttt{1.5} dans
  \texttt{producer/earthquake\_producer.py} et relancer le service
  (\texttt{docker compose restart earthquake-producer}).
  Avec M$\geq$1.5, un événement arrive environ toutes les 90 secondes.
  Pensez à remettre \texttt{2.5} une fois vos tests terminés.
\end{warningbox}

\subsection*{Questions --- Exploration de l'API}

\begin{enumerate}[label=\textbf{Q\arabic*.}]
  \item Combien de séismes de magnitude $\geq$ 5.0 se sont produits dans le monde
        au cours des 7 derniers jours ? (utilisez l'API directement)
  \item Dans le producer, comment convertit-on le timestamp USGS (en millisecondes epoch)
        en format ISO 8601 ? Quel module Python est utilisé ?
  \item Pourquoi le producer écrit-il dans \textbf{deux} destinations (Kafka ET PostgreSQL)
        plutôt que d'écrire uniquement dans Kafka et de laisser Flink peupler PostgreSQL ?
  \item La méthode \texttt{ON CONFLICT (event\_id) DO NOTHING} est utilisée lors
        de l'insertion PostgreSQL. Quel problème cela résout-il face au chevauchement
        de \texttt{POLL\_OVERLAP\_SEC} ?
\end{enumerate}

\newpage

% ===========================================================================
%  PARTIE 2 — DAG AIRFLOW
% ===========================================================================
\section{Construction du DAG Airflow}

\subsection{Rôle du DAG dans cette architecture}

Contrairement à une architecture batch classique, \textbf{Airflow n'ingère pas les données}
dans ce pipeline. L'ingestion est déléguée au service producer qui tourne en permanence.

Le DAG joue uniquement le rôle d'\textbf{initialisateur du pipeline} :

\begin{center}
\ttfamily\small
health\_check $\to$ create\_topics $\to$ run\_flink\_job $\to$ start\_tableau\_sink $\to$ verify\_output
\end{center}

\begin{conceptbox}{Séparation des responsabilités}
  \begin{itemize}
    \item \textbf{Producer (Docker, permanent)} : ingestion temps réel, 24h/24
    \item \textbf{DAG Airflow (déclenché une fois)} : orchestration du démarrage
  \end{itemize}
  Cette séparation est essentielle : Airflow est un orchestrateur, pas un moteur de traitement.
  Il ne doit pas bloquer un thread pendant 30 secondes pour attendre une API externe.
\end{conceptbox}

\subsection{Démarrage de l'infrastructure}

\begin{lstlisting}[style=bash, caption={Lancement des services Docker}]
cd infra/
docker compose up -d

# Attendre que tous les services soient healthy (~3 minutes)
docker compose ps
\end{lstlisting}

Vérifiez que ces services sont en état \texttt{healthy} ou \texttt{running} :
\texttt{kafka}, \texttt{flink-jobmanager}, \texttt{flink-taskmanager},
\texttt{flink-sql-client}, \texttt{airflow-webserver}, \texttt{airflow-scheduler},
\texttt{airflow-postgres}, \texttt{earthquake-producer}, \texttt{kafka-to-postgres}.

\begin{warningbox}
  Les services \texttt{earthquake-producer} et \texttt{kafka-to-postgres} démarrent
  \textbf{automatiquement} avec Docker Compose (\texttt{restart: unless-stopped}).
  Vous n'avez pas à les lancer manuellement.
  La tâche \texttt{start\_tableau\_sink} du DAG se contente de vérifier
  que \texttt{kafka-to-postgres} est bien en cours d'exécution.
\end{warningbox}

\subsection{Structure du fichier DAG à créer}

Créez le fichier \texttt{dags/earthquake\_pipeline.py}. Voici la structure attendue :

\begin{lstlisting}[style=python, caption={Squelette du DAG --- dags/earthquake\_pipeline.py}]
import os
from datetime import datetime, timedelta

from airflow import DAG
from airflow.operators.bash import BashOperator

# --- Configuration (lue depuis les variables d'environnement) ---
KAFKA_BOOTSTRAP = os.environ.get("KAFKA_BOOTSTRAP_SERVERS", "kafka:29092")

# --- default_args ---
default_args = {
    "owner": "etudiant",
    "depends_on_past": False,
    "retries": 2,
    "retry_delay": timedelta(minutes=1),
    "execution_timeout": timedelta(minutes=10),
}

with DAG(
    dag_id="earthquake_pipeline",
    default_args=default_args,
    description="Pipeline seismes : USGS API -> Kafka -> Flink SQL",
    start_date=datetime(2025, 1, 1),
    schedule_interval="*/30 * * * *",   # toutes les 30 minutes
    catchup=False,
    max_active_runs=1,
    tags=["earthquake", "kafka", "flink", "tp7"],
) as dag:

    health_check       = BashOperator(task_id="health_check",       bash_command="...")
    create_topics      = BashOperator(task_id="create_topics",      bash_command="...")
    run_flink_job      = BashOperator(task_id="run_flink_job",      bash_command="...")
    start_tableau_sink = BashOperator(task_id="start_tableau_sink", bash_command="...")
    verify_output      = BashOperator(task_id="verify_output",      bash_command="...")

    health_check >> create_topics >> run_flink_job >> start_tableau_sink >> verify_output
\end{lstlisting}

\subsection{Tâche 1 --- health\_check}

Cette tâche vérifie que Kafka et Flink répondent avant de démarrer.

\begin{lstlisting}[style=bash, caption={Commande bash pour health\_check}]
echo "=== Verification Kafka ==="
docker exec kafka kafka-topics \
    --bootstrap-server localhost:9092 --list
echo "=== Verification Flink ==="
curl -sf http://flink-jobmanager:8081/overview
echo "Health check OK"
\end{lstlisting}

\begin{conceptbox}{Pourquoi un health check ?}
  En orchestration, il est courant de valider que les services dépendants sont prêts
  avant d'exécuter des tâches qui en ont besoin. Si Kafka n'est pas encore disponible,
  la tâche \texttt{create\_topics} échouerait avec un message d'erreur peu clair.
  Un health check explicite donne un message d'erreur précis et permet un retry propre.
\end{conceptbox}

\begin{warningbox}
  \textbf{Réfléchissez :} que se passe-t-il si le health check réussit mais que Kafka
  tombe \emph{entre} cette tâche et \texttt{create\_topics} ?
  C'est pour cela que \texttt{retries=2} est configuré dans les \texttt{default\_args} :
  Airflow peut relancer automatiquement une tâche en échec.
  Combien de tentatives au total votre tâche a-t-elle le droit d'effectuer ?
\end{warningbox}

\subsection{Tâche 2 --- create\_topics}

Cette tâche crée les 3 topics Kafka nécessaires au pipeline.

\begin{lstlisting}[style=bash, caption={Création idempotente des topics}]
for TOPIC in earthquake-raw earthquake-enriched earthquake-alerts; do
    PARTITIONS=3
    if [ "$TOPIC" != "earthquake-raw" ]; then PARTITIONS=1; fi
    docker exec kafka kafka-topics \
        --bootstrap-server localhost:9092 \
        --create --if-not-exists \
        --topic "$TOPIC" \
        --partitions $PARTITIONS \
        --replication-factor 1
    echo "Topic $TOPIC pret."
done
\end{lstlisting}

\begin{conceptbox}{Idempotence et \texttt{--if-not-exists}}
  Le flag \texttt{--if-not-exists} garantit que la commande ne retourne pas d'erreur
  si le topic existe déjà. C'est la règle d'or en orchestration :
  \textbf{chaque tâche doit supporter les retries et les re-runs sans effet de bord}.
\end{conceptbox}

\begin{warningbox}
  \textbf{Testez par vous-même :} exécutez le DAG une première fois, puis une deuxième.
  Que se passe-t-il si vous oubliez le flag \texttt{--if-not-exists} ?
  Kafka renvoie-t-il une erreur ? La tâche échoue-t-elle ?
  Vérifiez le comportement dans Kafdrop (\url{http://localhost:9090}).
\end{warningbox}

\subsection{Tâche 3 --- run\_flink\_job}

Cette tâche soumet le fichier SQL au Flink SQL Client.

\begin{lstlisting}[style=bash, caption={Soumission du fichier SQL à Flink}]
set -eo pipefail
OUTPUT=$(docker exec flink-sql-client \
    bin/sql-client.sh -f /opt/flink/jobs/earthquake_processing.sql 2>&1)
echo "$OUTPUT"
if echo "$OUTPUT" | grep -q '\[ERROR\]'; then
    echo "Erreur detectee dans la sortie Flink SQL." >&2
    exit 1
fi
echo "Job Flink SQL soumis avec succes."
\end{lstlisting}

\begin{infobox}{Airflow orchestre, Flink traite}
  Airflow ne traite \textbf{pas} la donnée lui-même. Il délègue au cluster Flink
  via \texttt{docker exec flink-sql-client bin/sql-client.sh -f <fichier>}.
  Le fichier SQL est monté dans le conteneur via le volume \texttt{flink-jobs/}.
  Une fois soumis, Flink tourne indépendamment d'Airflow.
\end{infobox}

\begin{warningbox}
  \textbf{Problème à anticiper --- slots Flink épuisés.}

  Vous avez implémenté la soumission. Le DAG tourne toutes les 30 minutes.
  Ouvrez l'interface Flink (\url{http://localhost:8082}) après 3 ou 4 runs.
  Combien de jobs sont en état \texttt{RUNNING} ?

  \medskip
  Chaque exécution du DAG soumet un \textbf{nouveau} job sans arrêter le précédent.
  Au bout de quelques heures, tous vos task slots sont consommés et Flink refuse
  toute nouvelle soumission.

  \medskip
  \textbf{Réfléchissez :} avant de soumettre, comment savoir si le job tourne déjà ?
  Flink expose une API REST sur \texttt{http://flink-jobmanager:8081/jobs/overview}
  qui liste tous les jobs et leur état (\texttt{RUNNING}, \texttt{FINISHED}, \texttt{FAILED}...).

  \medskip
  \textbf{Autre piste :} pour identifier un job par son nom dans cette API,
  encore faut-il lui avoir donné un nom stable. Flink SQL permet de nommer
  un job via \texttt{SET 'pipeline.name' = 'mon-job'} \emph{avant} l'instruction
  \texttt{INSERT INTO}. Sans cela, Flink attribue un nom générique impossible à retrouver.
\end{warningbox}

\subsection{Tâche 4 --- start\_tableau\_sink}

Cette tâche vérifie que le service Docker \texttt{kafka-to-postgres} est bien en cours
d'exécution. Ce service démarre automatiquement avec \texttt{docker compose up -d} ---
la tâche joue uniquement le rôle d'un \textbf{health check}.

\begin{lstlisting}[style=bash, caption={Vérification du service kafka-to-postgres}]
docker ps --filter "name=^/kafka-to-postgres$" --filter "status=running" \
    | grep -q kafka-to-postgres \
    && echo "kafka-to-postgres service OK." \
    || (echo "kafka-to-postgres non demarre -- verifier docker compose." && exit 1)
\end{lstlisting}

\begin{conceptbox}{Pourquoi un service Docker plutôt qu'une tâche Airflow ?}
  \texttt{kafka-to-postgres} est une \textbf{boucle infinie} : il consomme Kafka en continu
  et insère dans PostgreSQL sans jamais s'arrêter.
  Une tâche Airflow doit \emph{se terminer} --- lancer une boucle infinie en tâche
  bloquerait le pipeline jusqu'au timeout de 10 minutes, puis la tâche échouerait.

  \medskip
  La bonne pratique : un démon appartient à Docker Compose (\texttt{restart: unless-stopped}).
  Airflow, lui, \textbf{orchestre} --- il ne gère pas des processus longue durée.
\end{conceptbox}

\begin{warningbox}
  \textbf{Attention aux accolades dans BashOperator.}
  Airflow utilise Jinja2 pour rendre les \texttt{bash\_command}.
  Les séquences \texttt{\{\{...\}\}} sont interprétées comme des variables Jinja2.
  \texttt{--format '\{\{.State.Running\}\}'} (syntaxe Go template de Docker) provoquerait
  une erreur Jinja2. Utilisez \texttt{docker ps --filter} à la place.
\end{warningbox}

\subsection{Tâche 5 --- verify\_output}

Cette tâche vérifie que le pipeline produit bien des données dans Kafka
et que PostgreSQL est peuplé par le producer.

\begin{warningbox}
  \textbf{Attention au timing.}
  Flink a besoin de quelques secondes après la soumission pour démarrer
  et commencer à traiter des messages. Si \texttt{verify\_output} s'exécute
  immédiatement après \texttt{run\_flink\_job}, il est possible qu'aucun message
  ne soit encore dans \texttt{earthquake-enriched}.

  \medskip
  Utilisez l'option \texttt{--timeout-ms} du consumer Kafka pour lui laisser
  le temps d'attendre des messages. Quelle valeur choisir ?
  Trop courte, la vérification échoue à tort (faux négatif).
  Trop longue, la tâche bloque inutilement.
\end{warningbox}

\begin{lstlisting}[style=bash, caption={Vérification de bout en bout}]
echo "=== 1/2 Verification Kafka earthquake-enriched ==="
OUTPUT=$(docker exec kafka kafka-console-consumer \
    --bootstrap-server localhost:9092 \
    --topic earthquake-enriched \
    --from-beginning \
    --max-messages 3 \
    --timeout-ms 30000 2>&1)
echo "$OUTPUT"
echo "$OUTPUT" | grep -q "Processed a total of [1-9]" || exit 1
echo "Kafka earthquake-enriched OK."

echo "=== 2/2 Verification PostgreSQL earthquake_events ==="
COUNT=$(docker exec airflow-postgres psql \
    -U airflow -d airflow -t -c \
    "SELECT COUNT(*) FROM earthquake_events;" 2>&1 | tr -d ' ')
echo "Lignes dans earthquake_events : $COUNT"
[ "$COUNT" -gt 0 ] || exit 1
echo "PostgreSQL OK - $COUNT seisme(s) disponible(s) pour Tableau."
\end{lstlisting}

\subsection{Exécution du DAG}

Une fois le fichier créé dans \texttt{dags/}, Airflow le détecte automatiquement.

\begin{lstlisting}[style=bash, caption={Déclenchement manuel du DAG}]
# Via le CLI Airflow
docker exec airflow-webserver airflow dags trigger earthquake_pipeline

# Ou depuis l'interface web : http://localhost:8080
# Login : admin / admin
\end{lstlisting}

\subsection*{Questions --- DAG Airflow}

\begin{enumerate}[label=\textbf{Q\arabic*.}]
  \item Le DAG est planifié avec \texttt{schedule\_interval="*/30 * * * *"}.
        Pourquoi ce choix est-il cohérent avec la fréquence du producer (\texttt{POLL\_INTERVAL = 30s}) ?
        En quoi la vérification d'idempotence dans \texttt{run\_flink\_job} est-elle
        indispensable avec ce schedule automatique ?
  \item Quelle est la différence entre \texttt{retries=2} dans \texttt{default\_args}
        et \texttt{execution\_timeout} ? Que se passe-t-il si Flink met
        12 minutes à soumettre le job ?
  \item La chaîne de dépendances est \texttt{health\_check >> create\_topics >> run\_flink\_job >> start\_tableau\_sink >> verify\_output}.
        Pourquoi ces tâches doivent-elles être \textbf{séquentielles} et non parallèles ?
  \item Dans \texttt{verify\_output}, on utilise \texttt{--timeout-ms 30000} pour le consumer Kafka.
        Que se passe-t-il si aucun message n'arrive dans les 30 secondes ?
        Est-ce un vrai échec de pipeline ou un faux négatif possible ? Expliquez.
\end{enumerate}

\subsection*{Exercice --- Ajout d'une tâche de monitoring}

Ajoutez une cinquième tâche \texttt{print\_summary} \textbf{après} \texttt{verify\_output}
qui affiche un résumé lisible dans les logs Airflow :

\begin{itemize}
  \item le nombre de séismes dans PostgreSQL
  \item le nombre de topics Kafka et leurs offsets courants
  \item l'état des jobs Flink (via \texttt{curl http://flink-jobmanager:8081/jobs})
\end{itemize}

\newpage

% ===========================================================================
%  PARTIE 3 — TRAITEMENT FLINK SQL
% ===========================================================================
\section{Traitement avec Flink SQL}

\subsection{Pourquoi enrichir sans fenêtre temporelle ?}

Les séismes sont des événements rares : globalement, un séisme de magnitude $\geq$ 2.5
se produit environ toutes les 15 à 30 minutes.

\begin{warningbox}
  Si vous utilisez une fenêtre \textbf{TUMBLE de 60 secondes}, la fenêtre se ferme
  après 60s sans recevoir de séisme et produit un résultat vide --- ou pire,
  ne produit rien du tout. Le topic de sortie restera silencieux pendant des minutes.

  \textbf{Solution} : enrichir chaque événement \emph{immédiatement} à l'arrivée,
  sans attendre la fermeture d'une fenêtre.
\end{warningbox}

\begin{conceptbox}{Streaming stateless vs Streaming avec fenêtre}
  \begin{itemize}
    \item \textbf{Avec fenêtre (TUMBLE/HOP)} : Flink accumule les messages pendant
          $\Delta$, puis émet un résultat agrégé à la fermeture.
          Utile pour : statistiques par période, comptages, moyennes temporelles.
    \item \textbf{Sans fenêtre (INSERT INTO ... SELECT sans GROUP BY)} :
          chaque message entrant génère \textbf{immédiatement} un message sortant.
          C'est un opérateur \textit{stateless} : latence quasi nulle.
          Utile pour : enrichissement, filtrage, transformation champ par champ.
  \end{itemize}
  Dans ce TP, nous utilisons l'approche \textbf{sans fenêtre} car les séismes
  sont rares et nous voulons voir chaque événement enrichi immédiatement.
\end{conceptbox}

\subsection{Connexion au SQL Client (optionnel)}

Si vous souhaitez tester des requêtes interactivement avant d'écrire le fichier SQL :

\begin{lstlisting}[style=bash, caption={Accès au Flink SQL Client interactif}]
docker exec -it flink-sql-client bin/sql-client.sh

# Le prompt Flink SQL apparait :
# Flink SQL>
\end{lstlisting}

\subsection{Table SOURCE --- earthquake\_raw}

Créez d'abord la table qui lit le topic Kafka \texttt{earthquake-raw}.
Elle contient les données brutes publiées par le producer.

\begin{lstlisting}[style=sql, caption={Table source --- earthquake\_raw}]
CREATE TABLE IF NOT EXISTS earthquake_raw (
    event_id     STRING,
    magnitude    DOUBLE,
    mag_type     STRING,
    place        STRING,
    event_time   STRING,
    latitude     DOUBLE,
    longitude    DOUBLE,
    depth_km     DOUBLE,
    significance INT,
    tsunami      INT,
    status       STRING,
    type         STRING,
    title        STRING,
    ingested_at  STRING,
    proc_time    AS PROCTIME()
) WITH (
    'connector'                    = 'kafka',
    'topic'                        = 'earthquake-raw',
    'properties.bootstrap.servers' = 'kafka:29092',
    'scan.startup.mode'            = 'earliest-offset',
    'format'                       = 'json'
);
\end{lstlisting}

\begin{conceptbox}{PROCTIME() --- Temps de traitement}
  \texttt{PROCTIME()} est une colonne virtuelle générée par Flink,
  représentant l'horloge système au moment où Flink traite le message.
  Elle avance toujours, indépendamment des données.
  Dans ce TP, elle est définie dans la table source mais
  \textbf{n'est pas utilisée dans les jobs sans fenêtre}.
  Elle serait utilisée si vous ajoutiez un job TUMBLE ultérieurement.
\end{conceptbox}

\subsection{Table SINK 1 --- earthquake\_enriched}

Créez la table de sortie pour les séismes enrichis.

\begin{lstlisting}[style=sql, caption={Table sink --- earthquake\_enriched}]
CREATE TABLE IF NOT EXISTS earthquake_enriched (
    event_id        STRING,
    magnitude       DOUBLE,
    mag_type        STRING,
    place           STRING,
    event_time      STRING,
    latitude        DOUBLE,
    longitude       DOUBLE,
    depth_km        DOUBLE,
    significance    INT,
    tsunami         INT,
    status          STRING,
    title           STRING,
    depth_category  STRING,
    severity        STRING,
    is_significant  INT,
    processed_at    STRING
) WITH (
    'connector'                    = 'kafka',
    'topic'                        = 'earthquake-enriched',
    'properties.bootstrap.servers' = 'kafka:29092',
    'format'                       = 'json'
);
\end{lstlisting}

\subsection{Table SINK 2 --- earthquake\_alerts}

\begin{lstlisting}[style=sql, caption={Table sink --- earthquake\_alerts}]
CREATE TABLE IF NOT EXISTS earthquake_alerts (
    event_id     STRING,
    event_time   STRING,
    magnitude    DOUBLE,
    place        STRING,
    depth_km     DOUBLE,
    tsunami      INT,
    severity     STRING,
    alert_time   STRING
) WITH (
    'connector'                    = 'kafka',
    'topic'                        = 'earthquake-alerts',
    'properties.bootstrap.servers' = 'kafka:29092',
    'format'                       = 'json'
);
\end{lstlisting}

\subsection{Job 1 --- Enrichissement temps réel (sans fenêtre)}

Ce job lit chaque événement de \texttt{earthquake-raw} et ajoute trois champs calculés.
Chaque séisme entrant produit \textbf{immédiatement} une ligne enrichie en sortie.

\begin{warningbox}
  \textbf{Nommez vos jobs.}
  Par défaut, Flink attribue un nom générique à chaque job soumis
  (\texttt{insert-into\_earthquake\_enriched\_1}, etc.).
  Si vous soumettez le fichier SQL plusieurs fois, vous ne pouvez pas distinguer
  les jobs entre eux dans l'interface ou l'API REST.

  \medskip
  Avant chaque instruction \texttt{INSERT INTO}, ajoutez :
  \begin{center}
    \texttt{SET 'pipeline.name' = 'nom-stable-du-job';}
  \end{center}
  Ce nom stable est ce qui permet à la tâche \texttt{run\_flink\_job} du DAG
  de vérifier si le job est déjà en cours d'exécution.
\end{warningbox}

\begin{lstlisting}[style=sql, caption={Job 1 --- enrichissement sans fenêtre}]
INSERT INTO earthquake_enriched
SELECT
    event_id,
    magnitude,
    mag_type,
    place,
    event_time,
    latitude,
    longitude,
    depth_km,
    significance,
    tsunami,
    status,
    title,

    -- Categorie de profondeur (classification sismologique standard)
    CASE
        WHEN depth_km < 70  THEN 'Superficiel (0-70 km)'
        WHEN depth_km < 300 THEN 'Intermédiaire (70-300 km)'
        ELSE                     'Profond (300+ km)'
    END AS depth_category,

    -- Niveau de sévérité
    CASE
        WHEN tsunami = 1      THEN 'TSUNAMI'
        WHEN magnitude >= 8.0 THEN 'Dévastateur'
        WHEN magnitude >= 7.0 THEN 'Majeur'
        WHEN magnitude >= 6.0 THEN 'Fort'
        WHEN magnitude >= 5.0 THEN 'Modéré'
        ELSE                       'Mineur'
    END AS severity,

    -- Seisme significatif selon le seuil USGS (sig > 500)
    CASE WHEN significance > 500 THEN 1 ELSE 0 END AS is_significant,

    CAST(CURRENT_TIMESTAMP AS STRING) AS processed_at

FROM earthquake_raw;
\end{lstlisting}

\subsection{Job 2 --- Alertes temps réel (filtre sans fenêtre)}

Ce job filtre les séismes dangereux. Seuls les séismes de magnitude $\geq$ 5.0
ou avec alerte tsunami sont transmis dans \texttt{earthquake-alerts}.

\begin{lstlisting}[style=sql, caption={Job 2 --- alertes avec filtre WHERE}]
INSERT INTO earthquake_alerts
SELECT
    event_id,
    event_time,
    magnitude,
    place,
    depth_km,
    tsunami,
    CASE
        WHEN tsunami = 1      THEN 'TSUNAMI'
        WHEN magnitude >= 8.0 THEN 'Dévastateur'
        WHEN magnitude >= 7.0 THEN 'Majeur'
        WHEN magnitude >= 6.0 THEN 'Fort'
        ELSE                       'Modéré'
    END AS severity,
    CAST(CURRENT_TIMESTAMP AS STRING) AS alert_time
FROM earthquake_raw
WHERE magnitude >= 5.0 OR tsunami = 1;
\end{lstlisting}

\begin{conceptbox}{INSERT stateless --- latence quasi nulle}
  Sans \texttt{GROUP BY} ni fenêtre temporelle, Flink ne maintient \textbf{aucun état}
  entre deux messages. Chaque message est traité indépendamment et produit immédiatement
  un résultat (si la clause \texttt{WHERE} est vraie pour le Job 2).
  La latence est typiquement de quelques millisecondes.

  Cette propriété simplifie aussi la reprise sur erreur : Flink peut reprendre
  depuis n'importe quel offset Kafka sans avoir à reconstruire un état accumulé.
\end{conceptbox}

\subsection{Assemblage dans le fichier SQL}

Le fichier \texttt{flink-jobs/earthquake\_processing.sql} doit contenir les blocs
dans cet ordre :

\begin{enumerate}
  \item \texttt{CREATE TABLE IF NOT EXISTS earthquake\_raw ...}
  \item \texttt{CREATE TABLE IF NOT EXISTS earthquake\_enriched ...}
  \item \texttt{CREATE TABLE IF NOT EXISTS earthquake\_alerts ...}
  \item (optionnel) \texttt{SET 'pipeline.name' = '...';}
  \item \texttt{INSERT INTO earthquake\_enriched SELECT ...}
  \item (optionnel) \texttt{SET 'pipeline.name' = '...';}
  \item \texttt{INSERT INTO earthquake\_alerts SELECT ... WHERE ...}
\end{enumerate}

Le flag \texttt{IF NOT EXISTS} garantit que le fichier peut être soumis plusieurs fois
sans erreur (idempotence).

\begin{warningbox}
  \textbf{Le piège du re-run.}
  Vous soumettez votre fichier SQL, tout fonctionne. Vous déclenchez le DAG une seconde fois.
  Ouvrez l'UI Flink (\url{http://localhost:8082}) : combien de jobs \texttt{RUNNING} voyez-vous ?

  \medskip
  Si vous en voyez 4 au lieu de 2, c'est que votre tâche \texttt{run\_flink\_job}
  n'est pas encore idempotente. Retournez à la section précédente et implémentez
  la vérification via l'API REST Flink avant toute soumission.
  Les tests unitaires du dossier \texttt{tests/} vous aideront à valider votre implémentation.
\end{warningbox}

\subsection{Vérification des jobs Flink}

\begin{lstlisting}[style=bash, caption={Vérification des jobs via l'API REST Flink}]
# Liste des jobs en cours d'execution
curl -s http://localhost:8082/jobs | python3 -m json.tool

# Interface web Flink : http://localhost:8082

# Consommer quelques messages depuis earthquake-enriched
docker exec kafka kafka-console-consumer \
    --bootstrap-server localhost:9092 \
    --topic earthquake-enriched \
    --from-beginning \
    --max-messages 3 \
    --timeout-ms 15000
\end{lstlisting}

\subsection*{Questions --- Flink SQL}

\begin{enumerate}[label=\textbf{Q\arabic*.}]
  \item Un séisme de magnitude 7.5 \textbf{avec} alerte tsunami arrive.
        Quelle valeur aura le champ \texttt{severity} dans \texttt{earthquake-enriched} ?
        Et dans \texttt{earthquake-alerts} ? Expliquez pourquoi les deux valeurs peuvent différer.
  \item Le job d'enrichissement n'a \textbf{pas} de clause \texttt{WHERE}.
        Cela signifie que tous les séismes, même de magnitude 2.5, passent dans
        \texttt{earthquake-enriched}. Quel avantage y a-t-il à garder tous les événements
        dans ce topic plutôt que de filtrer ?
  \item Dans la table \texttt{earthquake\_raw}, pourquoi déclare-t-on
        \texttt{scan.startup.mode = 'earliest-offset'} ?
        Que se passerait-il si on utilisait \texttt{'latest-offset'} à la place ?
  \item \texttt{CAST(CURRENT\_TIMESTAMP AS STRING)} génère l'heure de \textbf{traitement Flink},
        pas l'heure du séisme. Dans quel cas ces deux valeurs seraient-elles très différentes ?
        Quel champ du message source représente l'heure réelle du séisme ?
\end{enumerate}

\subsection*{Exercice --- Job 3) : statistiques par fenêtre}

Ajoutez un Job 3 qui calcule des statistiques globales par \textbf{fenêtre TUMBLE de 5 minutes}.

\begin{warningbox}
  Attention : avec des séismes rares ($\sim$1 toutes les 15--30 min), la plupart des
  fenêtres de 5 minutes seront \textbf{vides}. Ce job ne produira des résultats
  que lorsqu'un séisme arrive. C'est normal, et c'est pourquoi les Jobs 1 et 2
  (sans fenêtre) sont préférables pour un monitoring temps réel.
\end{warningbox}

Colonnes attendues dans la table sink \texttt{earthquake\_stats} :
\texttt{window\_start}, \texttt{window\_end}, \texttt{nb\_earthquakes},
\texttt{mag\_avg}, \texttt{mag\_max}, \texttt{mag\_min}, \texttt{nb\_tsunami}.

\newpage

% ===========================================================================
%  PARTIE 4 — MONITORING
% ===========================================================================
\section{Monitoring et vérification}

\subsection{Kafdrop --- Interface Kafka}

Kafdrop est une interface web légère pour inspecter les topics et messages Kafka.

\begin{infobox}{Accès à Kafdrop}
  \url{http://localhost:9090}
\end{infobox}

\textbf{Vérifications à effectuer :}
\begin{enumerate}
  \item Vérifiez que les 3 topics existent :
        \texttt{earthquake-raw}, \texttt{earthquake-enriched}, \texttt{earthquake-alerts}.
  \item Inspectez quelques messages dans \texttt{earthquake-raw} :
        les coordonnées et la magnitude sont-elles cohérentes ?
  \item Combien de partitions a le topic \texttt{earthquake-raw} (3 partitions) ?
        Pourquoi ce choix par rapport aux topics de sortie (1 partition) ?
  \item Inspectez un message dans \texttt{earthquake-enriched} :
        les champs \texttt{depth\_category}, \texttt{severity} et \texttt{is\_significant}
        sont-ils correctement calculés ?
\end{enumerate}

\subsection{Flink Web UI}

\begin{infobox}{Accès au Flink Dashboard}
  \url{http://localhost:8082}
\end{infobox}

\textbf{Vérifications à effectuer :}
\begin{enumerate}
  \item Vérifiez que les 2 jobs (enrichissement + alertes) sont en état \texttt{RUNNING}.
  \item Inspectez le graphe d'exécution du Job 1 : identifiez les opérateurs
        \textit{Source}, \textit{Calc}, \textit{Sink}.
  \item Vérifiez le compteur \textit{Records Sent} : est-il en train d'augmenter
        lorsque le producer publie de nouveaux séismes ?
\end{enumerate}

\subsection{Airflow Web UI}

\begin{infobox}{Accès à Airflow}
  \url{http://localhost:8080} --- Login : \texttt{admin} / \texttt{admin}
\end{infobox}

\textbf{Vérifications à effectuer :}
\begin{enumerate}
  \item Vérifiez que le DAG \texttt{earthquake\_pipeline} est visible dans la liste.
  \item Déclenchez le DAG manuellement et observez les tâches passer au vert une par une.
  \item Inspectez les logs de \texttt{verify\_output} :
        combien de séismes sont disponibles dans PostgreSQL ?
\end{enumerate}

\subsection{Vérification PostgreSQL}

\begin{lstlisting}[style=bash, caption={Requêtes de vérification PostgreSQL}]
# Connexion directe au conteneur PostgreSQL
docker exec -it airflow-postgres psql -U airflow -d airflow

# Compter les séismes
SELECT COUNT(*) FROM earthquake_events;

# Les 5 séismes les plus récents
SELECT event_id, magnitude, place, event_time
FROM earthquake_events
ORDER BY ingested_at DESC
LIMIT 5;

# Distribution par sévérité (à calculer en SQL)
SELECT
    CASE
        WHEN magnitude >= 7.0 THEN 'Majeur+'
        WHEN magnitude >= 5.0 THEN 'Modere'
        ELSE                       'Mineur'
    END AS categorie,
    COUNT(*) AS nb
FROM earthquake_events
GROUP BY 1
ORDER BY 2 DESC;
\end{lstlisting}

\newpage

% ===========================================================================
%  PARTIE 5 (BONUS) — TABLEAU
% ===========================================================================
\section{Visualisation avec Tableau Desktop}

\subsection{Connexion PostgreSQL $\to$ Tableau}

PostgreSQL sert de pont entre le pipeline temps réel et Tableau Desktop.
Le producer y écrit les séismes en continu. Tableau peut rafraîchir sa connexion
à tout moment pour voir les derniers événements.

\begin{infobox}{Paramètres de connexion Tableau}
  \begin{tabular}{ll}
    \textbf{Type}     & PostgreSQL \\
    \textbf{Serveur}  & \texttt{localhost} \\
    \textbf{Port}     & \texttt{5433} (port externe exposé par Docker) \\
    \textbf{Base}     & \texttt{earthquake} \\
    \textbf{Login}    & \texttt{airflow} \\
    \textbf{Mot de passe} & \texttt{airflow} \\
    \textbf{Tables}   & \texttt{earthquake\_alerts}, \texttt{earthquake\_stats} \\
  \end{tabular}
\end{infobox}

\begin{warningbox}
  Le port \textbf{5433} (et non 5432) est utilisé car le port 5432 est réservé
  à une éventuelle instance PostgreSQL locale sur votre machine.
  À l'intérieur du réseau Docker, les services communiquent sur le port 5432.
\end{warningbox}

\subsection{Dashboard proposé}

Créez un dashboard Tableau avec les visualisations suivantes :

\begin{enumerate}
  \item \textbf{Carte mondiale} : utilisez les champs \texttt{longitude} et \texttt{latitude}.
        Taille des points proportionnelle à la magnitude.
        Couleur selon le champ \texttt{tsunami} (1 = alerte rouge, 0 = bleu).

  \item \textbf{Distribution des magnitudes} : histogramme par tranches de 0.5
        (2.5--3.0, 3.0--3.5, etc.).

  \item \textbf{Séismes par région} : regroupez par les 20 premiers caractères
        du champ \texttt{place} pour obtenir une distribution géographique approximative.

  \item \textbf{Top 10 des séismes les plus forts} : tableau trié par magnitude décroissante
        avec les colonnes \texttt{magnitude}, \texttt{place}, \texttt{event\_time}, \texttt{depth\_km}.
\end{enumerate}

\subsection{Rafraîchissement des données}

Tableau Desktop ne rafraîchit pas automatiquement les données depuis PostgreSQL.
Pour voir les nouveaux séismes :

\begin{itemize}
  \item Cliquez sur \textbf{Données > Actualiser} dans Tableau.
  \item Ou utilisez \textbf{F5} sur la vue (selon la version).
\end{itemize}

\begin{conceptbox}{Pourquoi PostgreSQL plutôt qu'une connexion directe à Kafka ?}
  Tableau Desktop ne dispose pas de connecteur natif pour les topics Kafka.
  PostgreSQL sert de \textbf{matérialisation} du flux : au lieu de consommer
  directement Kafka (streaming), Tableau interroge une table SQL classique
  (batch). Cette approche est très courante en production :
  on appelle cela un \textit{Lambda Architecture} ou une \textit{table matérialisée}.
\end{conceptbox}

\subsection*{Questions --- Tableau}

\begin{enumerate}[label=\textbf{Q\arabic*.}]
  \item La base \texttt{earthquake} (peuplée par \texttt{kafka-to-postgres}) contient
        les tables \texttt{earthquake\_alerts} et \texttt{earthquake\_stats}.
        Ces tables viennent de Flink --- quels champs enrichis y trouvez-vous
        que vous ne trouveriez pas dans \texttt{earthquake\_events} (peuplée par le producer) ?
  \item Le service \texttt{kafka-to-postgres} utilise \texttt{auto\_offset\_reset='earliest'}.
        Que se passe-t-il si vous relancez la stack (\texttt{docker compose down -v}) ?
        Les données PostgreSQL sont-elles perdues ? Le service re-consomme-t-il depuis le début ?
  \item Quel est l'intervalle de rafraîchissement minimum raisonnable pour Tableau
        dans ce contexte ? Justifiez en vous appuyant sur la fréquence de publication
        du producer (\texttt{POLL\_INTERVAL = 30s}).
\end{enumerate}

\newpage

% ===========================================================================
%  PARTIE 5 — TESTS UNITAIRES
% ===========================================================================
\section{Tests unitaires du pipeline}

\subsection{Philosophie des tests}

Dans un pipeline de données temps réel, les tests ne peuvent pas dépendre
de Kafka, Flink ou PostgreSQL pour s'exécuter rapidement.
On distingue deux niveaux :

\begin{center}
\begin{tabularx}{\textwidth}{lX}
  \toprule
  \textbf{Niveau} & \textbf{Ce qu'il teste} \\
  \midrule
  \textbf{Tests unitaires} & Logique Python pure (parsing, calculs, retries) --- sans services externes \\
  \textbf{Tests d'intégration} & Flux bout en bout avec Kafka, Flink et PostgreSQL réels \\
  \bottomrule
\end{tabularx}
\end{center}

Dans ce TP, nous implémentons des \textbf{tests unitaires} dans \texttt{tests/}.

\subsection{Structure du dossier tests}

\begin{lstlisting}[style=bash, caption={Arborescence des tests}]
tests/
  __init__.py
  test_dag.py        # structure et configuration du DAG Airflow
  test_producer.py   # logique de parsing et de retry du producer
\end{lstlisting}

\subsection{Tests du DAG Airflow (\texttt{test\_dag.py})}

Ces tests utilisent le \texttt{DagBag} d'Airflow pour charger et inspecter le DAG
\textbf{sans exécuter aucune tâche}.

\begin{lstlisting}[style=python, caption={Exemple de test de structure du DAG}]
from airflow.models import DagBag

def test_dag_loads_without_errors():
    dagbag = DagBag(dag_folder="dags", include_examples=False)
    assert "dags/earthquake_pipeline.py" not in dagbag.import_errors

def test_dag_schedule(dag):
    assert dag.schedule_interval == "*/30 * * * *"

def test_linear_chain(dag):
    chain = ["health_check", "create_topics", "run_flink_job", "verify_output"]
    for i, task_id in enumerate(chain[:-1]):
        task = dag.get_task(task_id)
        downstream_ids = {t.task_id for t in task.downstream_list}
        assert chain[i + 1] in downstream_ids
\end{lstlisting}

\textbf{Tests couverts :}
\begin{itemize}
  \item Chargement du DAG sans erreur de parsing
  \item Schedule (\texttt{*/30 * * * *}), \texttt{catchup=False}, \texttt{max\_active\_runs=1}
  \item Présence et nombre exact des 4 tâches
  \item Chaîne de dépendances linéaire complète
  \item \texttt{retries=2} sur chaque tâche
  \item Paramètre \texttt{sql\_file} de la tâche \texttt{run\_flink\_job}
\end{itemize}

\subsection{Tests du producer (\texttt{test\_producer.py})}

Les dépendances \texttt{kafka}, \texttt{psycopg2} et \texttt{requests} sont
\textbf{mockées} avant l'import du module pour éviter toute dépendance réseau.

\begin{lstlisting}[style=python, caption={Mock des dépendances externes}]
import sys, types
from unittest.mock import MagicMock

kafka_mock = types.ModuleType("kafka")
kafka_mock.KafkaProducer = MagicMock
sys.modules.setdefault("kafka", kafka_mock)

import producer.earthquake_producer as ep
\end{lstlisting}

\textbf{Tests couverts :}
\begin{itemize}
  \item \texttt{\_parse\_feature} : coordonnées, timestamp ISO 8601, significance,
        fallback \texttt{code} $\to$ \texttt{id}, gestion des \texttt{None}
  \item \texttt{\_wait\_for\_kafka} : succès immédiat + \texttt{RuntimeError} après \texttt{max\_retries}
  \item \texttt{\_wait\_for\_postgres} : idem
\end{itemize}

\subsection{Exécution des tests}

\begin{lstlisting}[style=bash, caption={Lancement des tests avec uv}]
# Installer les dependances de test
uv pip install pytest apache-airflow==2.10.4

# Lancer tous les tests
uv run pytest tests/ -v

# Lancer uniquement les tests du DAG
uv run pytest tests/test_dag.py -v
\end{lstlisting}

\begin{warningbox}
  \texttt{uv run} utilise le Python défini dans \texttt{.python-version} (3.12).
  Si \texttt{apache-airflow} n'est pas installé dans le venv actif,
  \texttt{test\_dag.py} échoue avec \texttt{ModuleNotFoundError: No module named 'airflow'}.
  Vérifiez que votre venv est activé avant de lancer les tests.
\end{warningbox}

\subsection*{Questions --- Tests}

\begin{enumerate}[label=\textbf{Q\arabic*.}]
  \item Les tests du DAG n'exécutent aucune tâche Bash.
        Quelle limite cela implique-t-il ? Quel type de bug ne sera pas détecté
        par ces tests mais le serait par un test d'intégration ?
  \item Dans \texttt{test\_producer.py}, pourquoi mocke-t-on \texttt{kafka} et \texttt{psycopg2}
        \textbf{avant} d'importer le module producer ?
        Que se passerait-il si on les mockait après l'import ?
  \item La fonction \texttt{\_wait\_for\_kafka} effectue \texttt{max\_retries} tentatives
        avec un délai \texttt{delay} entre chaque. Dans les tests, on passe \texttt{delay=0}.
        Pourquoi est-ce important pour la rapidité des tests ?
\end{enumerate}

\newpage

% ===========================================================================
%  QUESTIONS DE RÉFLEXION
% ===========================================================================
\section*{Questions de réflexion}
\addcontentsline{toc}{section}{Questions de réflexion}

Ces questions n'ont pas de réponse unique --- elles visent à approfondir
votre compréhension des architectures de streaming.

\begin{enumerate}[label=\textbf{R\arabic*.}]

  \item \textbf{Temps réel vs orchestration batch.}
        Dans cette architecture, le producer tourne en permanence sans passer par Airflow.
        Si vous deviez utiliser Airflow pour l'ingestion (déclenché toutes les 15 minutes),
        quels problèmes cela poserait-il par rapport à un service continu ?
        Pensez aux notions de latence, de perte de données entre deux runs,
        et de charge sur l'API USGS.

  \item \textbf{Event Time vs Processing Time.}
        Un séisme peut être détecté par les sismographes et transmis à l'API
        avec un délai de quelques minutes.
        Le champ \texttt{processed\_at} calculé par Flink correspond au temps de traitement,
        pas au temps réel du séisme. Dans quel scénario cette différence serait-elle critique ?
        Comment la notion de \textit{watermark} et d'\textit{event time} permettrait-elle
        de corriger cela ?

  \item \textbf{Idempotence de bout en bout.}
        Le producer utilise \texttt{ON CONFLICT (event\_id) DO NOTHING} pour PostgreSQL.
        Flink, lui, peut re-traiter des messages depuis \texttt{earliest-offset} à chaque restart.
        Que se passe-t-il dans \texttt{earthquake-enriched} si Flink redémarre ?
        Y a-t-il des doublons ? Comment un consommateur Flink devrait-il se comporter
        pour être réellement idempotent ?

  \item \textbf{Passage à l'échelle.}
        Ce pipeline ingère aujourd'hui quelques dizaines de séismes par heure.
        Si vous deviez ingérer en temps réel 1 million d'événements par seconde
        (par exemple : tous les capteurs IoT d'un réseau sismique mondial),
        quels composants de l'architecture devriez-vous redimensionner ?
        Quels paramètres Kafka (partitions, replication-factor, retention)
        et Flink (parallelism, checkpointing) modifieriez-vous ?

\end{enumerate}

% ===========================================================================
%  ANNEXE — SCHÉMA DES TOPICS KAFKA
% ===========================================================================
\newpage
\appendix
\section{Schéma des messages Kafka}

\subsection{Topic : earthquake-raw}

Publié par le \textbf{producer} depuis l'API USGS.

\begin{center}
\begin{tabularx}{\textwidth}{llX}
  \toprule
  \textbf{Champ} & \textbf{Type} & \textbf{Description} \\
  \midrule
  \texttt{event\_id}   & STRING & Identifiant unique USGS \\
  \texttt{magnitude}   & DOUBLE & Magnitude du séisme \\
  \texttt{mag\_type}   & STRING & Type de magnitude (\texttt{ml}, \texttt{mw}, ...) \\
  \texttt{place}       & STRING & Description du lieu \\
  \texttt{event\_time} & STRING & Date/heure ISO 8601 UTC du séisme \\
  \texttt{latitude}    & DOUBLE & Latitude en degrés décimaux \\
  \texttt{longitude}   & DOUBLE & Longitude en degrés décimaux \\
  \texttt{depth\_km}   & DOUBLE & Profondeur du foyer en km \\
  \texttt{significance}& INT    & Indice USGS 0--1000+ \\
  \texttt{tsunami}     & INT    & 1 si alerte tsunami, 0 sinon \\
  \texttt{status}      & STRING & \texttt{automatic} ou \texttt{reviewed} \\
  \texttt{type}        & STRING & Type d'événement (ex: \texttt{earthquake}) \\
  \texttt{title}       & STRING & Titre complet USGS \\
  \texttt{ingested\_at}& STRING & Timestamp d'ingestion ISO 8601 UTC \\
  \bottomrule
\end{tabularx}
\end{center}

\subsection{Topic : earthquake-enriched}

Produit par \textbf{Flink Job 1} --- contient tous les champs de \texttt{earthquake-raw}
plus les champs calculés.

\begin{center}
\begin{tabularx}{\textwidth}{llX}
  \toprule
  \textbf{Champ} & \textbf{Type} & \textbf{Description} \\
  \midrule
  \multicolumn{3}{l}{\textit{(mêmes champs que earthquake-raw, sauf type et ingested\_at)}} \\
  \midrule
  \texttt{depth\_category} & STRING & \texttt{Superficiel (0-70 km)} / \texttt{Intermédiaire (70-300 km)} / \texttt{Profond (300+ km)} \\
  \texttt{severity}        & STRING & \texttt{TSUNAMI} / \texttt{Dévastateur} / \texttt{Majeur} / \texttt{Fort} / \texttt{Modéré} / \texttt{Mineur} \\
  \texttt{is\_significant} & INT    & 1 si \texttt{significance > 500}, sinon 0 \\
  \texttt{processed\_at}   & STRING & Timestamp de traitement Flink (ISO 8601) \\
  \bottomrule
\end{tabularx}
\end{center}

\subsection{Topic : earthquake-alerts}

Produit par \textbf{Flink Job 2} --- uniquement les séismes $\geq$ 5.0 ou avec tsunami.

\begin{center}
\begin{tabularx}{\textwidth}{llX}
  \toprule
  \textbf{Champ} & \textbf{Type} & \textbf{Description} \\
  \midrule
  \texttt{event\_id}   & STRING & Identifiant USGS \\
  \texttt{event\_time} & STRING & Heure du séisme (ISO 8601) \\
  \texttt{magnitude}   & DOUBLE & Magnitude \\
  \texttt{place}       & STRING & Lieu \\
  \texttt{depth\_km}   & DOUBLE & Profondeur \\
  \texttt{tsunami}     & INT    & 1 si alerte tsunami \\
  \texttt{severity}    & STRING & \texttt{TSUNAMI} / \texttt{Dévastateur} / \texttt{Majeur} / \texttt{Fort} / \texttt{Modéré} \\
  \texttt{alert\_time} & STRING & Heure d'émission de l'alerte (ISO 8601) \\
  \bottomrule
\end{tabularx}
\end{center}

\subsection{Table PostgreSQL : earthquake\_events}

Peuplée en temps réel par le \textbf{producer}. Utilisée par \textbf{Tableau Desktop}.

\begin{center}
\begin{tabularx}{\textwidth}{llX}
  \toprule
  \textbf{Colonne} & \textbf{Type} & \textbf{Description} \\
  \midrule
  \texttt{event\_id}   & TEXT PRIMARY KEY & Identifiant unique USGS \\
  \texttt{magnitude}   & DOUBLE PRECISION & Magnitude du séisme \\
  \texttt{mag\_type}   & TEXT & Type de magnitude \\
  \texttt{place}       & TEXT & Description du lieu \\
  \texttt{event\_time} & TEXT & Date/heure ISO 8601 UTC \\
  \texttt{latitude}    & DOUBLE PRECISION & Latitude \\
  \texttt{longitude}   & DOUBLE PRECISION & Longitude \\
  \texttt{depth\_km}   & DOUBLE PRECISION & Profondeur du foyer en km \\
  \texttt{significance}& INT & Indice USGS 0--1000+ \\
  \texttt{tsunami}     & INT & 1 si alerte tsunami \\
  \texttt{status}      & TEXT & \texttt{automatic} ou \texttt{reviewed} \\
  \texttt{type}        & TEXT & Type d'événement \\
  \texttt{title}       & TEXT & Titre complet USGS \\
  \texttt{ingested\_at}& TEXT & Timestamp d'ingestion ISO 8601 UTC \\
  \bottomrule
\end{tabularx}
\end{center}

\end{document}
